\documentclass[11pt]{article}
\title{CSP 203: Software Tools and Technologies \\
	(Question 1)}
\author{Bodike Chaithali (B24CS013)}
\date{\today}
\usepackage{amsmath,amssymb,amsthm,graphicx}
\usepackage{algorithmic}
\usepackage{float}
\begin{document}
	\maketitle
	\section{Introduction}
	Live streaming is a powerful and direct way to get connected with thou-
	sands of concurrent viewers. No matter if you are using it for your product
	promotion, giving a motivating speech, or just want to stream the game.
	\subsection{Usage}
	It is a tool to circulate information in the form of videos. Users can enjoy
	the audio and video without waiting for their file to download.
	\subsection{Statistics}
	As live is on rise Halder et al. \cite{size2020video}, the demand for high-quality and lower
	latency services is increasing. Thus, efficient, scalable, and economical solu-
	tions are required for live video streaming.
	
	
	According to the report of Business Fortune Insights, the global video
	streaming market size is expected to grow from USD 342.44 billion in 2019
	to USD 842.93 billion by 2027 at a compound annual growth rate of 12
	Size \cite{halder2021fybrrstream}.
	\section{Equation Section}
	The score of a node is calculated after taking uploading bandwidth, video
	streaming rate, latency, and stability of that node into consideration.
	\begin{enumerate}
		\item eqn. \eqref{eq:1} has been proposed for the calculation of score.
		\begin{equation}\label{eq:1}
			Score(peer)=k_1*(\frac{upload\_bandwidth}{streaming\_rate})+k_2*(\frac{1}{latency})+k_3*(\frac{active\_duration}{nums\_of\_failure})
		\end{equation}
	\item The general Tailor Series expansion is given in the form \eqref{eq:2}:
	\begin{equation}\label{eq:2}
		f(x)=f(a)+f'(a)(x-a)+\frac{f''(a)}{2!}(x-a)^2+\frac{f^{(3)}(a)}{3!}(x-a)^3+\cdots+\frac{f^{(n)}(a)}{n!}(x-a)^n+\cdots
	\end{equation}
	\item  The above expansion can be expressed as:
	\begin{equation*}
		f(x)=\sum_{n=0}^{\infty} \frac{f^{(n)}(a)}{n!}(x-a)^n
	\end{equation*}
	\end{enumerate}
	\section{Table}
	\begin{table}[h]
		\caption{Performance Comparison of Algorithms}
		\centering
		\begin{tabular}{|c|c|c|}
			\hline
			Algorithm & Time Complexity & Space Complexity \\ \hline
			Linear Search & $O(n)$ & $O(1)$ \\ \hline
			Binary Search & $O(\log n)$ & $O(1)$ \\ \hline
			Merge Sort & $O(n \log n)$ & $O(n)$ \\ \hline
			Quick Sort & $O(n \log n)$ & $O(\log n)$ \\ \hline
		\end{tabular}
	\end{table}
	\section{Algorithm}
	The following algorithm-\eqref{eq:2} is used to search an element in a sorted array
	in $O(\log(n))$ time.
	\section{Image}
	\begin{figure}[H]
		\label{fg:1}
		\centering
		\includegraphics[width=0.8\textwidth]{/home/user/Desktop/B24CS013_Bodike Chaithali_exam_03/B24CS013_Task1/B24CS013.png}
		\caption{Linux Desktop Image}
	\end{figure}


	The above fig.\ref{fg:1} can be seen after we log in to the ubuntu operating sys-
	tem.
	\section{Bibliography}
	\bibliography{B24CS013}
	\bibliographystyle{plain}
\end{document}